\section{Concluding Remarks}
\label{sec:concluding}

We have determined when the distinguished twisted signing minimizes
the adjacency spectral radius of the signed square of an even cycle.  The
answer is not monotone in the order: equality holds at every even order
$n\geq 8$ except $n=32$, $n=40$, and all even $n\geq 48$.  The proof also
identifies a structural mechanism behind the eventual failure.  A
period-eight reference phase admits an elementary six-gap interface whose
squared spectral edge is $c_6$, and separated copies of this interface can be
placed on large finite rings.  This mechanism proves the analytic tail, while
finite certificates determine the exceptional orders and the onset at
$48$.

The result leaves three natural questions.  First, along each nonzero even
residue class modulo $8$, does the sequence of minimum squared spectral radii
converge to $c_6$?  The present argument proves only the corresponding upper
limit bound, and supplies no matching lower bound.

Second, is the six-gap interface optimal in a broader class of
sector-changing finite cores?  Our comparison theorem is restricted to
single-gap interfaces, so it does not exclude a genuinely multi-gap core with
a smaller spectral edge.

Third, do spectral-radius minimizers at sufficiently large orders have an
asymptotic description in terms of well-separated elementary phase slips?
The constructions used here establish effective competitors, but they do not
classify the minimizing switching classes or force such a decomposition.
